\chapter{対角 Gaussian decoder と \texorpdfstring{$W_2$}{W2}}
\label{ch:gaussian-w2}

第\ref{ch:wasserstein-metrics}章では，一般の Polish 空間上で
$W_p$ が距離であることを示した．曲率へ進むには，二点間の距離だけでなく，
微小変位に対する二次形式が必要になる．そこで本章以降は
\[
  p=2,\qquad \mathcal{X}=\R^n
\]
に限定する．さらに，まず decoder 分布を対角 Gaussian に限定する．
この場合は $W_2$ が閉形式をもち，一般の Wasserstein 空間の
微分構造を準備せずに，潜在空間上の計量テンソルまで直接進めるからである．

\begin{remark}{Euclid 空間は Polish 空間である}{gw-euclidean-polish}
  以下 $\R^n$ には Euclid 距離 $d(x,y):=\norm{x-y}$ を入れる．
  $\mathbb{Q}^n$ は可算かつ $\R^n$ に稠密であり，Euclid 距離に関する
  Cauchy 列は各座標で収束するから $\R^n$ は完備である．
  よって定義~\ref{def:meas-polish}の意味で $\R^n$ は Polish 空間であり，
  第\ref{ch:wasserstein-metrics}章の $W_2$ の理論を適用できる．
\end{remark}

\section{対角 Gaussian 分布}
\label{sec:gw-gaussian}

\begin{definition}{標準 Gaussian 測度}{gw-standard-gaussian}
  $n\in\N$ とする．$\R^n$ 上の密度
  \[
    \varphi_n(x):=(2\pi)^{-n/2}\exp\left(-\frac{\norm{x}^2}{2}\right)
    \qquad(x\in\R^n)
  \]
  を Lebesgue 測度（定義~\ref{def:meas-lebesgue}）に関してもつ確率測度を
  \textbf{$n$ 次元標準 Gaussian 測度}といい，$\gamma^{(n)}$ と書く．
  すなわち，$A\in\Borel(\R^n)$ に対して
  \[
    \gamma^{(n)}(A)=\int_A\varphi_n(x)\,\d x
  \]
  である．
\end{definition}

\begin{remark}{標準 Gaussian の一次・二次モーメント}{gw-standard-moments}
  座標を $x=(x_1,\dots,x_n)$ とすると，
  \[
    \int_{\R^n}x_i\,\d\gamma^{(n)}(x)=0,
    \qquad
    \int_{\R^n}x_ix_j\,\d\gamma^{(n)}(x)
    =\begin{cases}1&(i=j),\\0&(i\neq j)\end{cases}
  \]
  が成り立つ．第一式と $i\neq j$ の場合は密度の各座標に関する対称性から，
  $i=j$ の場合は1次元 Gaussian 積分から従う．
\end{remark}

\begin{definition}{対角 Gaussian 分布}{gw-diagonal-gaussian}
  $m=(m_1,\dots,m_n)\in\R^n$，
  $\sigma=(\sigma_1,\dots,\sigma_n)\in(0,\infty)^n$ とする．
  アフィン写像
  \[
    A_{m,\sigma}(x)
    :=m+\operatorname{diag}(\sigma_1,\dots,\sigma_n)x
  \]
  による $\gamma^{(n)}$ の像測度
  \[
    \mathcal{N}\bigl(m,\operatorname{diag}(\sigma^2)\bigr)
    :=\gamma^{(n)}\circ A_{m,\sigma}^{-1}
  \]
  を，平均 $m$，対角共分散
  $\operatorname{diag}(\sigma_1^2,\dots,\sigma_n^2)$ の
  \textbf{対角 Gaussian 分布}という．
\end{definition}

\begin{lemma}{対角 Gaussian は有限二次モーメントをもつ}{gw-gaussian-p2}
  \[
    \mathcal{N}\bigl(m,\operatorname{diag}(\sigma^2)\bigr)
    \in\Pp{2}(\R^n)
  \]
  であり，
  \[
    \int_{\R^n}\norm{x}^2\,
    \d\mathcal{N}\bigl(m,\operatorname{diag}(\sigma^2)\bigr)(x)
    =\norm{m}^2+\norm{\sigma}^2
  \]
  が成り立つ．
\end{lemma}

\begin{proof}
  像測度の変数変換公式（補題~\ref{lem:meas-change-of-variables}）と
  注意~\ref{rem:gw-standard-moments}より，
  \[
    \begin{aligned}
      \int_{\R^n}\norm{x}^2\,
      \d\mathcal{N}\bigl(m,\operatorname{diag}(\sigma^2)\bigr)(x)
      &=\int_{\R^n}\sum_{i=1}^n(m_i+\sigma_i z_i)^2
        \,\d\gamma^{(n)}(z)\\
      &=\sum_{i=1}^n(m_i^2+\sigma_i^2)
       =\norm{m}^2+\norm{\sigma}^2<\infty.
    \end{aligned}
  \]
  基点を $0\in\R^n$ とする定義~\ref{def:wp-p-space}により結論を得る．
\end{proof}

\section{1次元 Gaussian 間の \texorpdfstring{$W_2$}{W2}}
\label{sec:gw-one-dimensional}

最初に1次元で計算する．上界を与える coupling を一つ作るだけでは
最適性は分からないため，任意の coupling に対する下界も同時に示す．

\begin{theorem}{1次元 Gaussian 間の $W_2$}{gw-w2-one-dimensional}
  $a,b\in\R$，$s,t>0$ とし，
  \[
    \mu=\mathcal{N}(a,s^2),
    \qquad
    \nu=\mathcal{N}(b,t^2)
  \]
  とおく．このとき
  \[
    \boxed{
      \Wass_2(\mu,\nu)^2=(a-b)^2+(s-t)^2
    }
  \]
  である．さらに，写像
  \[
    z\longmapsto(a+sz,\ b+tz)
  \]
  による $\gamma^{(1)}$ の像測度は最適 coupling である．
\end{theorem}

\begin{proof}
  \textbf{1. 任意の coupling に対する下界．}

  $\pi\in\Couplings(\mu,\nu)$ を任意に取る．周辺積分公式
  （補題~\ref{lem:meas-marginal-integral}）より
  \[
    \begin{gathered}
      \int(x-a)\,\d\pi(x,y)=0,
      \qquad
      \int(y-b)\,\d\pi(x,y)=0,\\
      \int(x-a)^2\,\d\pi(x,y)=s^2,
      \qquad
      \int(y-b)^2\,\d\pi(x,y)=t^2.
    \end{gathered}
  \]
  したがって
  \[
    \begin{aligned}
      \int(x-y)^2\,\d\pi(x,y)
      &=(a-b)^2+s^2+t^2
        -2\int(x-a)(y-b)\,\d\pi(x,y).
    \end{aligned}
  \]
  Hölder の不等式（定理~\ref{thm:meas-holder}，指数 $2,2$）より
  \[
    \int(x-a)(y-b)\,\d\pi(x,y)
    \leq
    \left(\int(x-a)^2\,\d\pi\right)^{1/2}
    \left(\int(y-b)^2\,\d\pi\right)^{1/2}
    =st.
  \]
  よってすべての $\pi\in\Couplings(\mu,\nu)$ に対して
  \[
    \int(x-y)^2\,\d\pi(x,y)
    \geq(a-b)^2+s^2+t^2-2st
    =(a-b)^2+(s-t)^2.            \tag{1}
  \]

  \textbf{2. 下界を達成する coupling．}

  $B(z):=(a+sz,b+tz)$ とおき，
  $\pi^*:=\gamma^{(1)}\circ B^{-1}$ とする．
  第1成分の周辺分布は $\mu$，第2成分の周辺分布は $\nu$ だから
  $\pi^*\in\Couplings(\mu,\nu)$ である．また，像測度の積分公式より
  \[
    \begin{aligned}
      \int(x-y)^2\,\d\pi^*(x,y)
      &=\int_{\R}\bigl((a-b)+(s-t)z\bigr)^2
        \,\d\gamma^{(1)}(z)\\
      &=(a-b)^2+(s-t)^2.
    \end{aligned}
  \]
  これは (1) の下界に一致するため，$\pi^*$ は最適であり，
  主張した $W_2^2$ の公式が得られる．
\end{proof}

\begin{remark}{最適 coupling の意味}{gw-common-noise}
  上の coupling は，同じ標準 Gaussian 変数 $z$ を用いて
  $x=a+sz$ と $y=b+tz$ を同時に作る．つまり二つの分布の
  同じ分位点どうしを対応させる coupling である．
  一般の最適 coupling は一意とは限らないが，$s,t>0$ の1次元
  Gaussian の場合は，Hölder の等号条件からこの coupling が一意になる．
  本章で必要なのは存在と閉形式であり，一意性は以下では用いない．
\end{remark}

\section{対角 Gaussian 間の閉形式}
\label{sec:gw-diagonal-formula}

\begin{theorem}{対角 Gaussian 間の $W_2$}{gw-w2-diagonal}
  $m,m'\in\R^n$，$\sigma,\sigma'\in(0,\infty)^n$ とし，
  \[
    \mu=\mathcal{N}\bigl(m,\operatorname{diag}(\sigma^2)\bigr),
    \qquad
    \nu=\mathcal{N}\bigl(m',\operatorname{diag}((\sigma')^2)\bigr)
  \]
  とおく．このとき
  \[
    \boxed{
      \Wass_2(\mu,\nu)^2
      =\norm{m-m'}^2+\norm{\sigma-\sigma'}^2
    }
  \]
  である．さらに，写像
  \[
    z\longmapsto
    \left(m+\operatorname{diag}(\sigma)z,\,
          m'+\operatorname{diag}(\sigma')z\right)
  \]
  による $\gamma^{(n)}$ の像測度は最適 coupling である．
\end{theorem}

\begin{proof}
  \textbf{1. 任意の coupling に対する下界．}

  $\pi\in\Couplings(\mu,\nu)$ を任意に取る．各 $i=1,\dots,n$ について，
  写像 $(x,y)\mapsto(x_i,y_i)$ による $\pi$ の像測度を $\pi_i$ とする．
  その周辺分布はそれぞれ
  $\mathcal{N}(m_i,\sigma_i^2)$ と
  $\mathcal{N}(m_i',(\sigma_i')^2)$ だから，
  $\pi_i$ はこの二つの1次元 Gaussian の coupling である．
  定理~\ref{thm:gw-w2-one-dimensional}の下界より
  \[
    \begin{aligned}
      \int_{\R^n\times\R^n}\norm{x-y}^2\,\d\pi(x,y)
      &=\sum_{i=1}^n\int_{\R\times\R}(x_i-y_i)^2\,\d\pi_i(x_i,y_i)\\
      &\geq\sum_{i=1}^n
        \bigl((m_i-m_i')^2+(\sigma_i-\sigma_i')^2\bigr)\\
      &=\norm{m-m'}^2+\norm{\sigma-\sigma'}^2.
    \end{aligned}                                      \tag{2}
  \]

  \textbf{2. 下界を達成する coupling．}

  主張中の写像を $C$ とし，
  $\pi^*:=\gamma^{(n)}\circ C^{-1}$ とおく．定義から
  $\pi^*\in\Couplings(\mu,\nu)$ であり，
  \[
    \begin{aligned}
      \int\norm{x-y}^2\,\d\pi^*(x,y)
      &=\int_{\R^n}
        \norm{(m-m')+\operatorname{diag}(\sigma-\sigma')z}^2
        \,\d\gamma^{(n)}(z)\\
      &=\norm{m-m'}^2+\norm{\sigma-\sigma'}^2.
    \end{aligned}
  \]
  これは (2) の下界を達成する．$W_2^2$ の定義から結論を得る．
\end{proof}

\begin{remark}{一般の Gaussian との違い}{gw-full-covariance}
  共分散行列が対角でない一般の Gaussian に対しても $W_2$ は閉形式をもつが，
  行列平方根を含む Bures--Wasserstein 距離が現れる．
  これは Givens--Shortt (1984) および Takatsu (2011) で扱われる．
  本資料では曲率計算までの最短経路を保つため，
  標準偏差を Euclid ベクトルとして扱える対角の場合に限定する．
\end{remark}

\section{Gaussian decoder が定める引き戻し距離}
\label{sec:gw-decoder-distance}

\begin{definition}{対角 Gaussian decoder}{gw-gaussian-decoder}
  潜在空間を $\mathcal{Z}\subset\R^d$ とする．二つの写像
  \[
    m\colon\mathcal{Z}\to\R^n,
    \qquad
    \sigma\colon\mathcal{Z}\to(0,\infty)^n
  \]
  に対して
  \[
    K_z:=\mathcal{N}\bigl(m(z),\operatorname{diag}(\sigma(z)^2)\bigr)
    \in\Pp{2}(\R^n)
  \]
  を潜在点 $z$ における\textbf{decoder 分布}という．また
  \[
    F\colon\mathcal{Z}\to\R^{2n},
    \qquad
    F(z):=(m(z),\sigma(z))
  \]
  を decoder の\textbf{パラメータ写像}という．
\end{definition}

\begin{proposition}{引き戻し距離の閉形式}{gw-pullback-distance}
  $z,z'\in\mathcal{Z}$ に対して
  \[
    d_K(z,z'):=\Wass_2(K_z,K_{z'})
  \]
  とおくと，
  \[
    \boxed{
      d_K(z,z')=\norm{F(z)-F(z')}
    }
  \]
  が成り立つ．したがって $d_K$ は $\mathcal{Z}$ 上の擬距離であり，
  $F$ が単射なら距離である．
\end{proposition}

\begin{proof}
  定理~\ref{thm:gw-w2-diagonal}を $K_z,K_{z'}$ に適用すると
  \[
    d_K(z,z')^2
    =\norm{m(z)-m(z')}^2+\norm{\sigma(z)-\sigma(z')}^2
    =\norm{F(z)-F(z')}^2.
  \]
  非負な平方根を取れば閉形式を得る．Euclid 距離の対称性と
  三角不等式から $d_K$ は擬距離である．さらに $F$ が単射なら
  \[
    d_K(z,z')=0
    \iff F(z)=F(z')
    \iff z=z'
  \]
  なので距離になる．これは一般の距離性
  （命題~\ref{prop:wp-metric}）をこの Gaussian 族上で具体化したものである．
\end{proof}

\begin{remark}{曲率へ進むために残る問題}{gw-transition-curvature}
  $d_K(z,z')$ は，$F(z),F(z')\in\R^{2n}$ を結ぶ直線の長さ，すなわち
  decoder 像の周囲の Euclid 空間で測った\textbf{弦長}である．
  曲率はこの二点間公式だけから計算するのではなく，
  $F$ の微分が各潜在方向の長さをどのように変えるかを表す
  局所的な計量テンソルから計算する．次章でこの微分を取り出す．
\end{remark}

\begin{remark}{文献との対応}{gw-references}
  $W_2$ の定義と距離性は Villani (2009) Definitions 6.1, 6.4 に従う．
  Gaussian 間の一般の閉形式は Givens--Shortt (1984) に遡り，
  Gaussian 測度の有限次元 Riemann 幾何は Takatsu (2011) で詳しく扱われる．
  本章の対角版は，後続の引き戻し計算に必要な範囲を
  定理~\ref{thm:gw-w2-one-dimensional}から自己完結に証明したものである．
\end{remark}
